% Relatório de até cinco páginas. Compilado por `python -m redes report` (src/redes/report.py),
% que grava `numeros.tex` (macros de números) ao lado deste arquivo antes de chamar o latexmk.
% Nenhum número de resultado é digitado aqui: use as macros de numeros.tex.
\documentclass[11pt,a4paper]{article}
\usepackage[T1]{fontenc}
\usepackage[utf8]{inputenc}
% O babel desta instalação não traz brazilian.ldf nem padrões de hifenização do português;
% os nomes em português vêm do arquivo de localização, e \emergencystretch evita linhas largas.
\usepackage[english]{babel}
\babelprovide[import=pt-BR,main]{brazilian}
\emergencystretch=3em
\usepackage{mathptmx}
\usepackage[scaled=.92]{helvet}
\usepackage[margin=2.3cm,top=2.4cm,bottom=2.3cm]{geometry}
\usepackage{microtype}
\usepackage{amsmath}
\usepackage{graphicx}
\usepackage{booktabs}
\usepackage{xcolor}
\usepackage{setspace}
\usepackage{fancyhdr}
\usepackage{titlesec}
% hyperref antes de abntex2cite: na ordem inversa o .aux quebra em \abntnextkey.
\usepackage[hidelinks]{hyperref}
\usepackage[alf,abnt-full-initials=yes,abnt-etal-text=it]{abntex2cite}

\input{numeros}

% Aparência do projeto: Times, seções em caixa alta e negrito, parágrafo recuado, número da
% página no canto superior direito.
\setstretch{1.08}
\setlength{\parindent}{1.0cm}
\setlength{\parskip}{0pt}
\pagestyle{fancy}
\fancyhf{}
\fancyhead[R]{\small\thepage}
\renewcommand{\headrulewidth}{0pt}
\fancypagestyle{plain}{\fancyhf{}\fancyhead[R]{\small\thepage}\renewcommand{\headrulewidth}{0pt}}
\titleformat{\section}{\normalfont\normalsize\bfseries\MakeUppercase}{\thesection}{0.6em}{}
\titlespacing*{\section}{0pt}{1.1ex plus .3ex}{0.6ex}
\newcommand{\paragrafo}[1]{\par\noindent\textbf{#1}\ }

% Legendas no padrão ABNT: título acima, fonte abaixo.
\newcounter{figura}
\newcommand{\legendafigura}[2]{\refstepcounter{figura}{\centering\small Figura~\thefigura\ --- #2\label{#1}\par}\vspace{0.8ex}}
\newcounter{quadro}
\newsavebox{\caixaquadro}
\newcommand{\legendaquadro}[2]{\refstepcounter{quadro}{\centering\small Quadro~\thequadro\ --- #2\label{#1}\par}\vspace{0.8ex}}
\newcommand{\fonteabnt}[1]{\par\vspace{0.6ex}{\centering\footnotesize #1\par}}

\newcommand{\figura}[2]{\includegraphics[width=#2]{#1}}

\begin{document}

\begin{center}
{\bfseries\MakeUppercase{Universidade de São Paulo}}\\
{\bfseries\MakeUppercase{Faculdade de Economia, Administração, Contabilidade e Atuária}}\\[0.5ex]
{\small EAE1429 --- Redes Econômicas e Sociais, 2.º semestre de 2026}\\[1.4ex]
{\large\bfseries\MakeUppercase{Redes empresariais e persistência econômica sob alternância política: evidências dos empenhos municipais no estado de São Paulo}}\\[1.4ex]
{\small\scshape Felipe Ferreira Alexandre --- 13963144 \quad Andrei Schoschlakow Bonifácio --- 14558482\\
Rafael Eiki Teruya --- 14604204 \quad João Pedro Ferreira --- 14582312}\\[0.4ex]
{\small São Paulo, 2026}\par
\avisoAmostra
\end{center}
\vspace{-0.6ex}\hrule\vspace{1.2ex}

\section{Motivação}
O acesso aos empenhos municipais sobrevive à alternância política? A troca de CNPJs pode
ocultar continuidade societária, compatível tanto com persistência de poder
\cite{acemoglu2008} quanto com capacidade produtiva ou reorganização empresarial. Estimamos o
efeito da derrota eleitoral sobre a alocação e examinamos a transmissão pela rede como hipótese.
Ampliamos a conexão da firma isolada \cite{boas2014,baltrunaite2020} para a rede de sócios,
em diálogo com estudos de elites e parentesco \cite{schisani2021,cruz2017}.
\sintesePrincipal

\section{Dados e rede}
\label{sec:dados}
Combinamos as despesas municipais do TCESP, as eleições e doações do TSE e o quadro de sócios da
Receita Federal \cite{tcesp2026,tse2026,rfb2026}. Os fornecedores anteriores são definidos pelos
empenhos de janeiro de 2021 a setembro de 2024. As doações entram pela data da receita, até 5 de
outubro, véspera da eleição. Pessoas jurídicas não doam desde a ADI 4650 \cite{adi4650}, então
as conexões passam por pessoas físicas. Para empresas sem sócios na fotografia de
setembro, a rede usa vínculos observados na extração de 12/10/2024, após a eleição, com entrada
até 13/09/2024. Esse complemento pode perder vínculos encerrados antes da extração.
São \empresasRede{} empresas, \pessoasRede{}
pessoas e \vinculosRede{} vínculos. O grafo bipartido inclui sócios PF/PJ e administradores;
duas empresas são vizinhas quando compartilham um desses vínculos. Uma empresa
tem conexão direta quando um sócio ou administrador PF doou ao candidato do grupo. Essa
ampliação operacional do vínculo societário literal é examinada na variante sem administradores. Definimos
$\delta_{ik}=\min_{j\in D_k}d_G(i,j)/2$, sendo $D_k$ as empresas diretamente conectadas e
$d_G$ a distância no bipartido; $\delta=0$ nas sementes e $\infty$ sem caminho (Figura~\ref{fig:exemplo}).

O cruzamento liga \taxaCruzamentoA{} dos doadores de A e \taxaCruzamentoB{} dos de B a empresas.
Em \munAnalisaveis{} municípios analisáveis, \munConexaoA{} têm conexão a A e \munConexaoB{} a B.
O maior componente reúne \fracGigante{} das
empresas; sem as pessoas cuja retirada divide componentes, \fracGiganteSemArt{}.
As distribuições de grau, componentes e distâncias estão em \texttt{rede\_*.csv}.
Na comparação com a fotografia recente filtrada pela data de entrada,
\fracVinculosPerdidos{} dos vínculos de setembro ficam de fora e \fracVinculosExtras{} dos
vínculos reconstruídos não aparecem em setembro. Essa reconstrução é apenas diagnóstica;
\fracVinculosLacuna{} dos vínculos presentes na reconstrução recente mas ausentes em setembro
também aparecem no preenchimento de outubro de
2024. A vinculação entre empenhos e ajustes cobre \taxaVinculacao{} do valor de 2025;
os desfechos usam todos os empenhos do painel, inclusive os não vinculados.

\begin{figure}[t]
\centering
\legendafigura{fig:exemplo}{Rede societária de exemplo}%
\figura{figura1_exemplo}{0.62\linewidth}
\fonteabnt{Fonte: elaboração própria. Círculos: pessoas; retângulos: empresas. Bruno liga X a Y; Carla liga Y a Z. Ana, sócia de X, doa a A; Daniel, sócio de Z, doa a B (setas tracejadas).}
\end{figure}

\section{Desenho empírico}
Usamos eleições municipais apertadas, desenho com precedente em \citeonline{oliveira2025}. A
decisão de doar pode beneficiar empresas vizinhas se a vantagem se transmitir e o doador não
apropriar todo o retorno; essa externalidade é uma previsão, não uma propriedade observada do
grafo \cite{jackson1996}. A variável de corte $M_m$ é a diferença de votos válidos,
em pontos percentuais, entre o candidato do grupo incumbente e o adversário mais votado; o
tratamento $T_m$ vale 1 quando o incumbente perde. Sob continuidade dos resultados potenciais em
$M_m = 0$, municípios com vitórias e derrotas apertadas são localmente comparáveis; o desenho
identifica o efeito local da derrota no corte, e não o efeito médio em disputas folgadas. O grupo
incumbente A segue o protocolo prévio (prefeito, partido, coligação), B é o adversário mais votado;
casos sem candidato identificável são excluídos e listados. Para resultado $Y$ e
categoria de conexão $\ell$,
\begin{equation}
  Y_{jm}^{\ell} = \alpha^{\ell} + \beta^{\ell} T_m + \gamma^{\ell} M_m + \delta^{\ell} T_m M_m + u_{jm}^{\ell},
  \qquad |M_m| \le b,
\end{equation}
estimada por regressão linear local, núcleo triangular, banda MSE-ótima $b$ e correção de viés
\cite{calonico2014}. As categorias usam IC robusto CCT; os contrastes reamostram municípios,
recalculando a correção com bandas de estimação e de viés fixas por categoria. Reportamos também
meia banda ($b/2$) e dobro ($2b$), o suporte em municípios, o teste de densidade
\cite{cattaneo2020}, placebos com a participação pré-eleitoral e o efeito mínimo detectável.
H1 compara fornecedores anteriores diretamente conectados ao incumbente com os sem caminho
a ambos os grupos ($\delta_{iA}=\delta_{iB}=\infty$); H2
compara os efeitos por faixa de $\delta$. As participações de 2025 usam o empenho líquido
positivo do elemento. Em mercados sem empenho positivo, os fornecedores anteriores ficam na
amostra com participação zero. H3 decompõe a variação do empenho da rede dividida pelo gasto anual
médio pré-eleitoral do município: $Y_{Rm}=(E_{Rm,2025}-\bar E_{Rm,pre})/\bar E_{m,pre}$,
com barras indicando médias anuais. H4 examina a manutenção da posição das principais redes. As comparações
secundárias recebem correção de Holm.

%%RESULTADOS-INICIO
\section{Resultados}
O Quadro~\ref{tab:estimativas} reúne os contrastes. Em H2, as coortes têm distância fixa a cada
grupo. A sobreposição em mercado e porte
das coortes de conexão direta é \sobreposicaoDireta{}. O contraste
$(\hat\beta_B(d)-\hat\beta_A(d))/2$ resume heterogeneidade eleitoral entre coortes, que podem
compartilhar firmas e diferir na ligação ao outro grupo; não identifica o efeito de reduzir a distância.
A análise de sensibilidade usa apenas estratos de município, elemento e porte presentes nas
duas coortes (\texttt{h2\_comparaveis.csv}). \leituraHDoisComparaveis{}
Em H3, comparamos redes ligadas aos dois grupos (AMBOS) com as demais. O total soma fornecedoras
ligadas a A (R1), demais fornecedoras anteriores (R2) e entrantes (R3). Em H4, a comparação
descritiva em disputas de até cinco p.p. agrupa redes com e sem ligação ao vencedor.
No limiar central, as taxas de permanência correspondentes são \permanenciaComLigacao{} e
\permanenciaSemLigacao{}. No corte, o contraste usa as redes de ligação exclusiva, cujo estado
muda com o resultado, usando a mesma meia diferença B--A. Redes principais cobrem o limiar do
gasto acumulado e fornecem em três dos quatro anos; em 2025 o ranking inclui entrantes.
As taxas descritivas são médias em grupos distintos, sem interpretação causal.
\tabelaEstimativas
\tabelaCategorias
\leituraResultados

\begin{figure}[t]
\centering
\legendafigura{fig:rdd}{Contrastes entre coortes de distância e decomposição da variação da rede}%
\figura{gradiente_h2}{0.38\linewidth}\hfill\figura{decomposicao_h3}{0.38\linewidth}
\fonteabnt{Fonte: elaboração própria. Participação no elemento em H2; fração do gasto anual médio pré-eleitoral em H3. Barras indicam IC de 95\% por reamostragem de municípios.}
\end{figure}

\section{Robustez}
\tabelaRobustez
\leituraRobustez
\leituraSensibilidades
%%RESULTADOS-FIM

\section{Discussão e limitações}
\sinteseDiscussao{}
Os resultados são locais aos municípios paulistas do TCESP, sem a capital. Doações omitem
conexões informais; compartilhar sócios não mede controle. Homônimos geram erro de cruzamento,
e empenho não é contrato. Participações medem realocação: ganhos de uma firma reduzem a fatia
das demais. O agrupamento municipal absorve dependência interna, mas não interferência entre
municípios da mesma rede. Poucas disputas limitam o poder; perdas semelhantes também podem
refletir desorganização administrativa \cite{akhtari2022}. Só divulgamos agregados pseudonimizados.
Em H4, bandas amplas e ICs além do suporte probabilístico revelam instabilidade da extrapolação
corrigida; a análise permanece exploratória.

\bibliographystyle{abntex2-alf}
\begingroup\small\setstretch{1}
\bibliography{referencias}
\endgroup

\end{document}
